% !PRINTED 517
\weberpage{517}
Wenn $\omega$ prim ist, so ist der Satz evident, weil $\omega$ durch sich selbst theilbar ist. Wenn aber $\omega$ nicht prim ist, so ist es durch ein ganzes Functional $\omega_1$ theilbar, das weder eine Einheit, noch mit $\omega$ associirt ist. Ist also
\weq{1}{\omega = \omega_1 \alpha,}
so ist weder $\omega_1$ noch $\alpha$ eine Einheit. Aus (1) folgt aber
\[
N_a(\omega) = N_a(\omega_1)\, N_a(\beta), % [sic] printed with $\beta$, though (1) introduces $\alpha$
\]
und da $N_a(\beta)$ grösser als 1 ist, so ist
\weq{2}{N_a(\omega_1) < N_a(\omega).}

Wenn $\omega_1$ noch nicht prim ist, so kann man denselben Schluss auf $\omega_1$ anwenden und findet einen Theiler $\omega_2$ von $\omega_1$ derart, dass
\weq{3}{N_a(\omega_2) < N_a(\omega_1) < N_a(\omega)}
ist. Da es aber nur eine endliche Anzahl natürlicher Zahlen giebt, die kleiner sind als $N_a(\omega)$, so muss die Reihe der Functionale $\omega, \omega_1, \omega_2, \ldots$ abbrechen, und dies ist nur möglich, wenn das letzte von ihnen ein Primfunctional ist, wodurch der Satz 1. bewiesen ist.

2. \gesperrt{Jedes ganze von Null und von den Einheiten verschiedene Functional $\omega$ im Körper $\Omega$ kann in eine endliche Anzahl von Primfactoren zerlegt werden.}

Ist nämlich $\pi_1$ ein Primfactor von $\omega$ und
\weq{4}{\omega = \pi_1 \omega_1,}
so ist, da $N_a(\pi_1) > 1$ ist,
\[
N_a(\omega) > N_a(\omega_1).
\]

Ist $\omega_1$ keine Einheit, so ist es durch ein Primfunctional $\pi_2$ theilbar, und aus
\[
\omega_1 = \pi_2 \omega_2
\]
folgt
\[
N_a(\omega_1) > N_a(\omega_2).
\]

Fährt man so fort, so erhält man eine Reihe ganzer Functionale $\omega_1, \omega_2, \ldots$, deren absolute Normen fortwährend abnehmen, und diese Reihe bricht also mit einer Einheit ab. Ist $\pi_\nu$ die letzte von ihnen, die keine Einheit ist, so ist $\pi_\nu$ selbst ein Primfunctional, und es folgt
\weq{5}{\omega = \pi_1 \pi_2 \ldots \pi_\nu,}
w. z. b. w.
